\chapter{Registro das hipóteses}\label{ap:registro}

As hipóteses e os critérios de decisão foram registrados em quatro etapas, de 6 a 8 de novembro de 2026, cada uma antes do cálculo do que ela acrescentou. O texto de cada registro foi gravado com um resumo SHA-256, e uma mudança só pode entrar como nova emenda, sem alterar o registro anterior. Nenhum critério mudou depois dos resultados.

O registro de 6 de novembro contém H1 a H3, com os critérios do Quadro~\ref{tab:hipoteses} e o procedimento das seções~\ref{sec:descritores} a~\ref{sec:injecao}, e foi feito depois das análises exploratórias de 2 e 5 de novembro (seção~\ref{sec:preliminares}). Além dos critérios, ele definiu análises descritivas, que não entram na decisão. Em H1, elas são o sentido inverso, com treino na segunda rodada e teste na primeira, e o detector secundário. O registro já apontava que a primeira rodada da guilhotina, com \campoJanelasGuilhotinaUm\ janelas, sustenta um limiar estimado com poucos escores. Em H2, essa rodada, com menos de \num{25} janelas, entra só como descrição, e a menor mudança detectável também é expressa em porcentagem da mediana de $d_1$. Em H3, as análises descritivas são a assinatura no dobro da rotação, somada ao eixo axial, o detector secundário e o Isolation Forest com dez sementes diferentes do gerador de números aleatórios.

A primeira emenda, de 7 de novembro, acrescentou H4 e as análises posteriores AP1 e AP2. Quando ela foi escrita, eram conhecidos os resultados de H1 a H3, a mediana de $d_5$ em cada ensaio, o intervalo entre janelas consecutivas e, da base pública, só o artigo e os cadernos dos autores, sem nenhum dado. Para H4, a emenda manteve os descritores, o detector e o limiar da coleta, com o pré-processamento da seção~\ref{sec:base-publica}. As análises descritivas previstas para H4 são o menor nível de desbalanceamento com metade das janelas marcadas, a correlação de Spearman entre o fator de desbalanceamento e a fração marcada, o detector secundário, o Isolation Forest com dez sementes diferentes, o uso de um único canal de vibração, o adotado pelos autores da base, e a banda completa, até \SI{1000}{\hertz}, para medir o efeito do corte em \SI{184}{\hertz}. Os procedimentos da AP1 e da AP2 estão na seção~\ref{sec:res-posteriores}.

A segunda emenda, do mesmo dia, trocou a fonte dos dados de H4, porque a versão 151.2 da base deixou de estar disponível. Em 7 de novembro de 2026, o repositório da Fraunhofer-Gesellschaft estava fora do ar, e a cópia dessa versão no Internet Archive está incompleta. A versão 151.3, de dezembro de 2022, tem cópia completa no Internet Archive, e os arquivos comuns às duas versões que puderam ser comparados têm o mesmo tamanho e o mesmo código de verificação. Quando a emenda foi registrada, eram conhecidos os resultados de H1 a H3, da AP1 e da AP2, e a obtenção da versão 151.3 já tinha começado, sem nenhum descritor calculado. A hipótese, os critérios, o pré-processamento e as análises descritivas de H4 não mudaram.

A terceira emenda, de 8 de novembro, acrescentou quatro análises posteriores, AP3 a AP6, pedidas por pareceres sobre uma versão anterior deste texto, sem hipótese nova e sem mudança de critério. Quando ela foi registrada, eram conhecidos os resultados de H1 a H4, da AP1 e da AP2, o falso alarme de H4 com as sementes 0 a 9, de \hQuatroFaSementeMin~\% a \hQuatroFaSementeMax~\%, e a mediana de $d_1$ com o sensor em repouso. A AP3 recalcula o falso alarme de H1 e de H4 com as sementes 0 a 99 e com o treino reamostrado por blocos, refazendo o limiar e o modelo em cada réplica. A AP4 dá o intervalo da menor mudança detectável de H2, e a AP5 calcula a menor mudança detectável com o sensor em repouso, pela mesma regra de H2. A AP6 mede o escore do Isolation Forest em pontos fora da faixa do treino e o compara com o limiar. Os vereditos de H1 a H4, calculados com a semente 0, continuam os registrados.
